\documentclass[11pt,a4paper]{article}
\usepackage{shared/luxcover}
\usepackage[utf8]{inputenc}
\usepackage[T1]{fontenc}
\usepackage{amsmath,amssymb,amsfonts,amsthm}
\usepackage{graphicx}
\usepackage{hyperref}
\usepackage{booktabs}
\usepackage{enumitem}
\usepackage{xcolor}

\hypersetup{colorlinks=true,linkcolor=black,citecolor=blue,urlcolor=blue}

\newtheorem{theorem}{Theorem}[section]
\newtheorem{definition}{Definition}[section]
\newtheorem{remark}{Remark}[section]

\title{LP-309: Quantum Network Comparison\\[4pt]
\large Per-Chain Quantum Vulnerability Analysis and the\\
Case for Post-Quantum Infrastructure}

\author{
Zach Kelling\thanks{Corresponding author: zach@lux.network}\\
\textit{Lux Industries, Inc.}\\
\texttt{research@lux.network}
}

\date{Original: January 2025\quad|\quad Revised: March 2026}

\begin{document}

\luxcoverpage

\begin{abstract}
We analyze the quantum vulnerability of four major blockchain
networks---Bitcoin, Ethereum, Solana, and Cardano---examining their
key generation, address derivation, and transaction signing processes
to determine the window of public key exposure available to a quantum
adversary. We then describe how Lux Network's architecture mitigates
these vulnerabilities through post-quantum consensus signatures,
threshold key management with continuous rotation, and an X-Chain UTXO
model that minimizes key exposure.

This is not a hypothetical exercise. NIST estimates that
cryptanalytically relevant quantum computers may emerge within 10--15
years~\cite{nist-pqc-report}. Assets on-chain today will still be
on-chain then. The time to migrate is now.
\end{abstract}

\tableofcontents
\newpage

%% ============================================================
\section{Introduction}

Every major blockchain network relies on elliptic curve cryptography
(ECDSA over secp256k1, or EdDSA over Ed25519) for transaction signing.
Shor's algorithm~\cite{shor1999} can solve the elliptic curve discrete
logarithm problem in polynomial time on a fault-tolerant quantum
computer, enabling an adversary to derive private keys from exposed
public keys.

The critical question is not \emph{whether} this is possible, but
\emph{how long} each network exposes public keys to attack. This
exposure window determines the practical risk: a network where public
keys are visible for seconds has a fundamentally different threat
profile than one where they are visible indefinitely.

%% ============================================================
\section{Per-Chain Analysis}

\subsection{Bitcoin}

\subsubsection{Key Generation and Address Derivation}

Bitcoin generates an ECDSA key pair $(sk, pk)$ on secp256k1. The
address is derived through two hash layers:

\begin{enumerate}[nosep]
  \item $h_1 = \text{SHA-256}(pk)$
  \item $h_{160} = \text{RIPEMD-160}(h_1)$
  \item $\text{address} = \text{Base58Check}(\text{version} \| h_{160}
        \| \text{checksum})$
\end{enumerate}

\subsubsection{Public Key Exposure}

The public key $pk$ is \emph{not} revealed when receiving Bitcoin. The
address is a hash of $pk$, and hash preimage resistance protects $pk$
even if the address is known. The public key is revealed only when
\emph{spending}---the transaction includes $pk$ for signature
verification.

\subsubsection{Quantum Vulnerability Window}

The exposure window is the time between transaction broadcast and
block confirmation ($\sim$10 minutes for one confirmation, $\sim$60
minutes for six). An adversary must derive $sk$ from $pk$ within this
window to redirect funds.

\begin{remark}
Bitcoin's recommended practice of single-use addresses provides
additional protection: once an address's public key is revealed by
spending, no further funds should be sent to it. Unspent outputs at
never-used addresses have zero public key exposure.
\end{remark}

\subsection{Ethereum}

\subsubsection{Key Generation and Address Derivation}

Ethereum generates an ECDSA key pair on secp256k1. The address is
the last 20 bytes of the Keccak-256 hash of the public key:

\[
\text{address} = \text{Keccak-256}(pk)[12{:}32]
\]

\subsubsection{Public Key Exposure}

The public key is \textbf{effectively always exposed}. It can be
recovered from any ECDSA signature via ecrecover:

\[
pk = \text{ecrecover}(\text{msgHash}, v, r, s)
\]

Once an account has signed any transaction, the public key is
permanently derivable from the blockchain history.

\subsubsection{Quantum Vulnerability Window}

The exposure window is \textbf{permanent} for any account that has
ever sent a transaction. An adversary does not need to race against
block confirmation---they can extract $pk$ from historical transactions
at leisure and compute $sk$ offline.

\begin{remark}
Ethereum's account model exacerbates the risk. Unlike Bitcoin's UTXO
model where each output can use a fresh key, Ethereum accounts are
long-lived and typically reuse the same key for all transactions.
\end{remark}

\subsection{Solana}

\subsubsection{Key Generation and Address Derivation}

Solana uses Ed25519 key pairs. Unlike Bitcoin (SHA-256 + RIPEMD-160)
or Ethereum (Keccak-256), Solana applies \emph{no hash layer at all}:
the address is the raw public key, base58-encoded.

\[
\text{address} = \text{base58}(pk)
\]

\subsubsection{Public Key Exposure}

The public key is \emph{always} visible to anyone who knows the
address---there is no hash-based protection, even before first spend.
This is strictly weaker than Bitcoin's UTXO model (pk hidden until
spend) and equivalent in outcome to Ethereum's permanent exposure.

\subsubsection{Quantum Vulnerability Window}

Permanent from the moment an address is published. A quantum adversary
with Shor's algorithm can recover the secret key from the address
directly, without needing a prior transaction. Solana has the
\emph{worst} public-key exposure profile of the four chains analysed
here.

\subsection{Cardano}

\subsubsection{Key Generation and Address Derivation}

Cardano uses separate spending and staking key pairs. Addresses are
derived via Blake2b-224:

\begin{align*}
\text{PaymentKeyHash} &= \text{Blake2b-224}(\text{PublicSpendingKey}) \\
\text{StakeKeyHash} &= \text{Blake2b-224}(\text{PublicStakingKey}) \\
\text{BaseAddress} &= \text{Bech32}(\text{NetworkID} \|
  \text{PaymentKeyHash} \| \text{StakeKeyHash})
\end{align*}

\subsubsection{Public Key Exposure}

The spending public key is revealed when signing a transaction.
Cardano's extended UTXO model allows address reuse, but best practice
is to derive fresh addresses for each transaction.

\subsubsection{Quantum Vulnerability Window}

Similar to Bitcoin: the public key is exposed at signing time. However,
Cardano's faster block times ($\sim$20 seconds) reduce the exposure
window compared to Bitcoin's 10 minutes.

%% ============================================================
\section{Comparative Summary}

\begin{center}
\begin{tabular}{lcccc}
\toprule
\textbf{Property} & \textbf{BTC} & \textbf{ETH} & \textbf{SOL} &
  \textbf{ADA} \\
\midrule
Signature scheme & ECDSA & ECDSA & Ed25519 & Ed25519 \\
Address hash & SHA-256 + & Keccak-256 & none & Blake2b-224 \\
 & RIPEMD-160 & & (base58 pk) & \\
TX model & UTXO & Account & Account & Extended UTXO \\
PK exposure at receive & No & No & \textbf{Yes} & No \\
PK exposure at spend & Yes & Yes & Yes & Yes \\
PK recoverable from history & At spend & Always & Always & At spend \\
Exposure window & Minutes & Permanent & Permanent (from receive) & Seconds \\
Recommended key reuse & No & Yes (forced) & Yes (forced) & No \\
\midrule
\textbf{PQ resistance} & \textbf{Partial} & \textbf{Low} &
  \textbf{Worst} & \textbf{Partial} \\
\bottomrule
\end{tabular}
\end{center}

\subsection{Key Observations}

\begin{enumerate}[nosep]
  \item \textbf{UTXO models are inherently more resistant} than account
        models because they naturally support key-per-output patterns.
        The public key can be kept hidden until the moment of spending.
  \item \textbf{Account models force permanent exposure}: Once a key
        is used, it is recoverable forever. Migrating to a new key
        requires moving all assets, which is itself a transaction that
        exposes the old key.
  \item \textbf{Hash function security is not the bottleneck}: SHA-256,
        Keccak-256, and Blake2b-224 are all believed to be
        quantum-resistant (Grover provides only a quadratic
        speedup~\cite{grover1996}, so 256-bit hashes provide $\sim$128
        bits of quantum security). The vulnerability is in the
        \emph{signature scheme}, not the address derivation.
  \item \textbf{The real target is frozen funds}: Addresses with known
        public keys and large balances (exchange cold wallets, long-term
        holders, lost key situations) are the highest-value quantum
        targets. These funds sit on-chain with their public keys
        exposed indefinitely.
\end{enumerate}

%% ============================================================
\section{Lux Network's Post-Quantum Architecture}

\subsection{Consensus Layer}

Quasar consensus (LP-105) uses a hybrid signature scheme:

\begin{itemize}[nosep]
  \item \textbf{BLS-381}: Aggregate signatures for validator voting.
        Efficient ($\sim$48 bytes per signature), enabling compact block
        headers. Vulnerable to quantum attack via discrete log on
        pairing-friendly curves.
  \item \textbf{ML-DSA-65} (FIPS 204): NIST-standardized lattice-based
        signatures. 3.3\,KB signatures. Quantum-resistant under LWE
        hardness assumption.
  \item \textbf{Corona}: Lattice-based ring signatures providing
        signer anonymity within a validator set. Post-quantum by
        construction.
\end{itemize}

The hybrid design provides defense in depth: if BLS is broken by a
quantum adversary, ML-DSA and Corona signatures on the same blocks
prevent consensus forgery.

\subsection{X-Chain: UTXO with Minimal Exposure}

Lux's X-Chain inherits the UTXO model. Each transaction output is
associated with a fresh key. Public keys are revealed only at spending
time, and the network's sub-second finality (via Quasar) minimizes
the exposure window to under one second.

\begin{remark}
Compared to Bitcoin's 10-minute exposure window, Lux's sub-second
finality reduces the quantum attack window by approximately
$600\times$. This is not a substitute for post-quantum signatures---it
is an additional defense layer.
\end{remark}

\subsection{EVM Precompiles}

The Lux EVM provides post-quantum cryptographic operations as
precompiled contracts in the reserved \texttt{0x0200} range (see
\texttt{lux-evm-precompile-engineering.tex} for the canonical
specification, addresses, and gas costs). Available primitives
include ML-DSA-65 / SLH-DSA-128s / ML-KEM-768 verify and decapsulate,
FROST and CGGMP21 threshold verification, Corona ring verification,
and TFHE evaluation gates.

Smart contracts can verify post-quantum signatures natively, without
off-chain computation or trusted oracles.

\subsection{Bridge: MPC with Key Rotation}

The Lux Bridge uses 2-of-3 MPC threshold signing for cross-chain
transactions. Key rotation occurs at every withdrawal event:

\begin{enumerate}[nosep]
  \item Validators perform distributed key generation (DKG) for a new
        signing key.
  \item The withdrawal transaction spends from the old multisig address
        and sends remaining funds to a new address derived from the new
        key.
  \item The old key's public component is exposed for the duration of
        the withdrawal transaction only.
\end{enumerate}

This ensures no ECDSA public key is exposed for more than the time
required to confirm the withdrawal transaction. Combined with a
policy of batching withdrawals (hourly or daily), the total key
exposure across a year is measured in hours, not the permanent
exposure of Ethereum accounts.

\subsection{Hybrid Encryption: X-Wing}

For key encapsulation and encrypted communication, Lux uses X-Wing
(X25519 + ML-KEM-768 hybrid) implemented in \texttt{luxfi/age}. This
provides post-quantum confidentiality: an adversary must break
\emph{both} the elliptic curve and lattice components to decrypt.

%% ============================================================
\section{Applications Beyond Finance}

Post-quantum security is critical across multiple sectors where data
encrypted today must remain confidential for decades:

\begin{itemize}[nosep]
  \item \textbf{Government}: Classified document authentication,
        secure inter-agency communication, e-voting integrity.
  \item \textbf{Defense}: Secure command channels, supply chain
        verification, satellite communication authentication.
  \item \textbf{Healthcare}: Long-term patient record confidentiality,
        genomic data protection, cross-institution research without
        data exposure (via FHE).
  \item \textbf{Finance}: Transaction signing, key management for
        custodial assets, regulatory reporting with privacy guarantees.
\end{itemize}

The ``harvest now, decrypt later'' attack is not theoretical---it is
the baseline assumption for any adversary with long-term strategic
interests. Data encrypted under RSA or ECDSA today can be stored and
decrypted when quantum computers become available.

%% ============================================================
\section{Conclusion}

The analysis demonstrates a clear hierarchy of quantum vulnerability
among existing blockchain networks:

\begin{enumerate}[nosep]
  \item \textbf{Lowest risk}: UTXO-based networks with single-use
        addresses (Bitcoin, Lux X-Chain). Public key exposure is
        transient and can be minimized by key hygiene.
  \item \textbf{Highest risk}: Account-based networks (Ethereum,
        Solana). Public keys are permanently recoverable from
        transaction history.
\end{enumerate}

Neither category is \emph{safe}---both rely on ECDSA or EdDSA,
which fall to Shor's algorithm. The defense requires replacing the
signature scheme itself, not merely minimizing exposure windows.

Lux Network's approach is the only production architecture that
addresses all three layers: post-quantum \emph{consensus} signatures
(ML-DSA + Corona), post-quantum \emph{key encapsulation} (X-Wing),
and post-quantum \emph{smart contract verification} (EVM precompiles).
The UTXO model and sub-second finality provide additional defense in
depth, but they are not substitutes for correct cryptography.

\bibliographystyle{plain}
\begin{thebibliography}{99}

\bibitem{shor1999}
P.~W.~Shor, ``Polynomial-time algorithms for prime factorization and
discrete logarithms on a quantum computer,'' \textit{SIAM Review},
vol.~41, no.~2, pp.~303--332, 1999.

\bibitem{grover1996}
L.~K.~Grover, ``A fast quantum mechanical algorithm for database
search,'' \textit{STOC}, pp.~212--219, 1996.

\bibitem{nist-pqc-report}
NIST, ``Status Report on the Third Round of the NIST Post-Quantum
Cryptography Standardization Process,'' NISTIR~8413, 2022.

\bibitem{nist-pqc}
NIST, ``Post-Quantum Cryptography Standardization,''
\url{https://csrc.nist.gov/projects/post-quantum-cryptography}, 2024.

\bibitem{bernstein2017}
D.~J.~Bernstein and T.~Lange, ``Post-quantum cryptography---dealing
with the fallout of physics success,'' \textit{IACR Cryptology ePrint
Archive}, 2017/314.

\bibitem{regev2005}
O.~Regev, ``On lattices, learning with errors, random linear codes,
and cryptography,'' \textit{STOC}, pp.~84--93, 2005.

\bibitem{micciancio2009}
D.~Micciancio and O.~Regev, ``Lattice-based cryptography,'' in
\textit{Post-Quantum Cryptography}, pp.~147--191, Springer, 2009.

\bibitem{lyubashevsky2013}
V.~Lyubashevsky, C.~Peikert, and O.~Regev, ``On ideal lattices and
learning with errors over rings,'' \textit{Journal of the ACM},
vol.~60, no.~6, pp.~1--43, 2013.

\bibitem{lp105}
Z.~Kelling, ``LP-105: Quasar---Hybrid BLS + Lattice Consensus for
Post-Quantum Finality,'' Lux Industries, 2026.

\bibitem{lp305}
Z.~Kelling, ``LP-305: Quantum Secure Assets---A Post-Quantum
Methodology for Protecting Digital Assets on the Lux Network,''
Lux Industries, 2022--2026.

\end{thebibliography}

\end{document}
